Each step gives the setting we used and what it changed in our comparison (Table~\ref{tab:key}).
\begin{enumerate}\setlength{\itemsep}{1pt}
  \item \textbf{Classification: choose the head on the frozen encoder first.} Train a linear and an attentive probe and keep the better one on selection data: the attentive head was slightly better on TAP-CT ($\numCiClsattnlinTapctMeandiff$) and SPECTRE ($\numCiClsattnlinSpectreMeandiff$), the linear head clearly better on VoCo ($\numCiClsattnlinVocoMeandiff$ for the attentive head).
  \item \textbf{Classification: then adapt the attention weights under that head.} Rank-16 LoRA ($\alpha=32$, learning rate $10^{-4}$) on every block added $\numCiClsloralinAllMeandiff$ AUROC over the linear probe and $\numCiClsloraattnAllMeandiff$ over the attentive probe, positive in all \numCiClsloraattnAllClustersN{} cohort--backbone pairs; on VoCo, adapt under the linear head (Lin$\to$LoRA, $\numCiClslinthenloraglinVocoMeandiff$). Fine-tuning the attention weights in full was ahead of LoRA in \numSeedsclsAttnftLoraAllPosN{} of \numSeedsclsAttnftLoraAllPairsN{} seed pairs, by $\numCiClsattnftloraAllMeandiff$ on average, at about fifty times the trainable parameters. Use plain LoRA: no variant improved on it.
  \item \textbf{Classification: train one model per cohort, not one per label.} It was on average level with per-label models ($\numCiClsjointAllMeandiff$) at $1/K$ of the training.
  \item \textbf{Segmentation: put a multi-depth decoder on the frozen encoder.} It added $\numSegdeltaSpectreLinDecMeandiff$ to $\numCiSegdeclinVocoMeandiff$ Dice over a linear probe on every dataset. Adapting the encoder beneath it changed Dice by at most $\numCiSegloradecSpectreMeandiff$ on average; full fine-tuning ranked best but lost more than 0.1 Dice on some small structures, so check Dice per structure.
  \item \textbf{Segmentation: train a joint model with LoRA and the steps of the per-structure models it replaces.} At matched compute it was ahead of per-structure LoRA by $\numCmTotalsegSpectreLoraJointdiff$ to $\numCmAmosBBSpectreLoraJointdiff$ Dice; at our shorter budgets one model per structure was more accurate.
  \item \textbf{Full fine-tuning: lower the learning rate of early blocks.} With the encoder learning rate of the other methods it did not train on SPECTRE (near-constant predictions); layer-wise decay of 0.75 per block fixed it. Check that held-out predictions vary between cases.
\end{enumerate}

\paragraph{Limitations.} Each method has one untuned default setting, and only full fine-tuning was retrained after it failed, so the ranking among the stronger methods may change with tuning; seeds were repeated only for joint models. Segmentation had not converged and is scored on each backbone's working grid, so it ranks adaptation methods rather than measuring the distance to nnU-Net \citep{isensee2021nnunet} or to models pretrained on segmentation labels \citep{li2024suprem,he2024vista3d}. SPECTRE's and VoCo's pretraining images include several of our evaluation sets (Table~\ref{tab:overlap}), and one author co-developed TAP-CT.
